\begin{lemma}
Let \(G\) be a finite graph and let \(\nu\) be a probability law on
activation--priority pairs \((A,\pi)\).  Define
\[
I(A,\pi)=\{v\in A:\pi(v)>\pi(w)\text{ for every }w\in A\cap N_G(v)\}.
\]
Then the push-forward law \(\mu(S)=\nu\{(A,\pi):I(A,\pi)=S\}\) is a
probability mass function on subsets of \(V(G)\), satisfies the exact
push-forward identity for every \(S\subseteq V(G)\), and gives positive mass
only to independent sets of \(G\).
\end{lemma}
\begin{proof}
The set of all subsets \(S\subseteq V(G)\) is finite.  The events
\[
E_S=\{(A,\pi):I(A,\pi)=S\}
\]
are pairwise disjoint and their union is the whole sample space, because every
sample has a unique output set \(I(A,\pi)\).  Define \(\mu(S)=\nu(E_S)\),
viewed as a real number.  Since \(\nu\) is a probability law, every
\(\mu(S)\) is nonnegative, and finite additivity over the disjoint partition
\((E_S)_S\) gives
\[
\sum_S \mu(S)=\nu\left(\bigcup_S E_S\right)=\nu(\Omega)=1.
\]
The displayed definition is exactly the push-forward identity.

Now suppose \(\mu(S)>0\).  Then \(E_S\) is nonempty up to positive measure, so
there are samples with output \(S=I(A,\pi)\).  For any such sample, two
adjacent vertices cannot both lie in \(I(A,\pi)\).  Indeed, if \(u\) and \(v\)
are adjacent and both are in \(I(A,\pi)\), then the survival condition for
\(u\) gives \(\pi(v)<\pi(u)\), while the survival condition for \(v\) gives
\(\pi(u)<\pi(v)\), a contradiction.  Thus every output set is independent, and
hence every set receiving positive \(\mu\)-mass is independent.
\end{proof}
